\chapter{Project Management and Team Contribution}
\label{chap:project_management}

\section{Introduction}
The project was organized as a sequence of requirements, data, design, implementation, evaluation, and reporting milestones. Work was divided between Padmaja Shinde and Akshata Naik while key integration and validation tasks were completed jointly.

\section{Work Distribution}
\begin{table}[H]
\centering
\caption{Review 3 Work Distribution}
\begin{tabular}{|p{4cm}|p{4.5cm}|p{4.5cm}|}
\hline
\textbf{Area} & \textbf{Padmaja Shinde} & \textbf{Akshata Naik} \\
\hline
Core pipeline & PDF processing, preprocessing, chunking, retrieval, integration & Functional verification and bilingual workflow validation \\
\hline
Language support & Translation integration and intent handling & Translation fidelity and English/Hindi query testing \\
\hline
Quality & Backend debugging and source-grounding tests & Output, page evidence, classification, and unsupported-query tests \\
\hline
Documentation & Technical implementation content & Test evidence, results review, and presentation support \\
\hline
\end{tabular}
\end{table}

\section{Project Timeline and Milestones}
\begin{longtable}{|p{3cm}|p{5cm}|p{5cm}|}
\hline
\textbf{Phase} & \textbf{Milestone} & \textbf{Deliverable} \\
\hline
Review 1 & Problem definition, literature, dataset planning & Chapters 1--5 and proposed architecture \\
\hline
Review 2 & Prototype implementation and experimentation & Chapters 6--8, application modules, preliminary tests \\
\hline
Review 3 & Evaluation, error analysis, final results, impact, conclusion & Chapters 9--14, evidence appendices, final PDF and source archive \\
\hline
\end{longtable}

\section{Team Coordination}
The team used task-based ownership, shared review of interfaces between modules, and joint end-to-end testing. Findings from testing were converted into implementation changes and regression cases so that fixes could be verified after later edits.

\section{Testing and Quality Assurance}
Quality assurance includes automated unit and integration tests, an eight-document bilingual holdout, source-excerpt checks, value-preservation checks, and manual inspection of representative outputs. The complete automated suite was rerun for Review 3 and passed 18/18 tests.

\section{Version / Progress Tracking}
Source code and report material are maintained as version-controlled project files. Progress is tracked through implementation modules, test scripts, generated evaluation samples, and successive Review 1, Review 2, and Review 3 report chapters. Sensitive configuration values are separated from source through environment settings.

\section{Individual Contribution Summary}
Padmaja's primary contribution is the core document and NLP processing pipeline. Akshata's primary contribution is bilingual interaction, validation, testing, and system-level verification. Both members contributed to debugging, integration, documentation, and the final evaluation narrative. Detailed evidence is provided in Chapter 8 and Appendix E.

\section*{Chapter Summary}
The staged plan and complementary team roles enabled the project to progress from design to a tested final prototype. Traceable tests and artifacts provide evidence for the reported contributions.
